\chapter{1977年天津卷}

\section{解答题}

\begin{problem}
\begin{enumerate}
\item 在什么条件下，\(\frac{y}{2x}\)

\begin{circlelist}
\item 是正数；
\item 是负数；
\item 等于零；
\item 没有意义？
\end{circlelist}

\item 比较下列各组数的大小，并说明理由.

\begin{circlelist}
\item \(\cos31^\circ\) 与 \(\cos30^\circ\).
\item \(\log_2 1\) 与 \(\log_2 \frac{1}{4}\).
\end{circlelist}

\item 求值：

\begin{circlelist}
\item \(\tan\left(5\arcsin\frac{\sqrt{3}}{2}\right)\).
\item \(\left(-2\right)^0\times\left(0.01\right)^{\frac{1}{2}}\).
\end{circlelist}

\item 计算：\(\lg 12.5-\lg\frac{5}{8}+\lg\sin 30^\circ\).

\item 解方程：\(\frac{4x}{x^2-4}-\frac{2}{x-2}=1-\frac{1}{x+2}\).
\end{enumerate}
\end{problem}

\begin{solution}
\begin{enumerate}
\item 分式的分母为 \(2x\)，所以必须先有 \(x\ne0\).

当 \(x\ne0\) 时，\(2>0\)，故 \(\frac{y}{2x}>0\Longleftrightarrow xy>0\)，即 \(x\)，\(y\) 同号；同理，\(\frac{y}{2x}<0\Longleftrightarrow xy<0\)，即 \(x\)，\(y\) 异号.

当 \(y=0\) 且 \(x\ne0\) 时，\(\dfrac{y}{2x}=0\). 当 \(x=0\) 时分母为零，分式没有意义.

\item 在区间 \([0^\circ,180^\circ]\) 上，余弦函数随角度增大而减小. 因为 \(31^\circ>30^\circ\)，所以
\(\cos31^\circ<\cos30^\circ\).

又因为 \(\log_2 1=0\)，\(\log_2\frac{1}{4}=-2\)，所以 \(\log_2 1>\log_2\dfrac{1}{4}\).

\item 由反正弦函数的值域可知
\(\arcsin\frac{\sqrt{3}}{2}=60^\circ\).
因此
\[
\tan\left(5\arcsin\frac{\sqrt{3}}{2}\right)
=\tan300^\circ=-\sqrt{3}
\]

又因为 \((-2)\ne0\)，所以 \((-2)^0=1\)，且
\(\left(0.01\right)^{\frac{1}{2}}=0.1\)，
故
\(\left(-2\right)^0\left(0.01\right)^{\frac{1}{2}}=\frac{1}{10}\).

\item 因为 \(\sin30^\circ=\dfrac{1}{2}\)，由对数的运算法则，
\[
\lg 12.5-\lg\frac{5}{8}+\lg\sin30^\circ
=\lg\left(12.5\cdot\frac{\frac{1}{2}}{\frac{5}{8}}\right)=\lg10=1
\]

\item 原方程的定义域为 \(x\ne2\) 且 \(x\ne-2\). 因为
\(x^2-4=\left(x-2\right)\left(x+2\right)\)，
所以在定义域内，原方程左边为
\[
\frac{4x}{\left(x-2\right)\left(x+2\right)}-\frac{2}{x-2}
=\frac{4x-2\left(x+2\right)}{\left(x-2\right)\left(x+2\right)}
=\frac{2}{x+2}
\]
于是 \(\frac{2}{x+2}=1-\frac{1}{x+2}\).
由于 \(x+2\ne0\)，两边同乘 \(x+2\)，得 \(2=x+1\)，所以 \(x=1\). 代入原方程可知成立，故方程的解为 \(x=1\).
\end{enumerate}
\end{solution}

\begin{problem}
\begin{enumerate}
\item 某工厂准备在仓库的一侧建立一个矩形储料场（如图），现有 \(50\text{ 米}\) 长的铁丝网，如果用它来围成这个储料场，那么长和宽各是多少时，这个储料场的面积最大？并求出这个最大的面积.

\item 如图，已知 \(AB\)，\(DE\) 是圆 \(O\) 的直径，\(AC\) 是弦，\(AC\parallel DE\)，求证：\(CE=EB\).

\end{enumerate}
\bitmapfigure[width=4.6cm]{img/1977/tianjin/q2_1_fig1.png}

\bitmapfigure[width=6.0cm]{img/1977/tianjin/q2_2_fig1.png}
\end{problem}

\begin{solution}
\begin{enumerate}
\item 设储料场垂直于仓库的一边为 \(x\)，平行于仓库的一边为 \(y\). 因为仓库的一边可以直接作为围栏的一边，所以铁丝网只需围三边. 由围栏长度关系 \(2x+y=50\)，得 \(y=50-2x\).
由于 \(x>0\) 且 \(y>0\)，有 \(0<x<25\). 储料场面积为
\[
S=xy=x\left(50-2x\right)=-2x^2+50x
=-2\left(x-\frac{25}{2}\right)^2+\frac{625}{2}
\]
所以当 \(x=\dfrac{25}{2}\mathrm{~m}\) 时，面积取得最大值. 此时 \(y=50-2\cdot\frac{25}{2}=25\mathrm{~m}\)，且 \(S_{\max}=\frac{625}{2}\mathrm{~m}^2\).
即长为 \(25\mathrm{~m}\)，宽为 \(\dfrac{25}{2}\mathrm{~m}\) 时，面积最大.

\item 建立直角坐标系，使圆心 \(O\) 为原点，直径 \(DE\) 在 \(y\) 轴上，设圆的半径为 \(R\)，并取 \(E=\left(0,R\right)\). 设
\(A=\left(u,v\right)\).
由于 \(AB\) 是直径，\(O\) 是 \(AB\) 的中点，所以
\(B=\left(-u,-v\right)\).
又因为 \(AC\parallel DE\)，而 \(DE\) 在 \(y\) 轴上，所以 \(AC\) 是竖直弦，点 \(C\) 与点 \(A\) 的横坐标相同. 圆的方程为 \(x^2+y^2=R^2\)，同一条竖直直线与圆的两个交点的纵坐标互为相反数，故
\(C=\left(u,-v\right)\).
于是由两点间距离公式，\(CE^2=u^2+\left(-v-R\right)^2\)，\(EB^2=\left(-u\right)^2+\left(-v-R\right)^2\).
两式右端相等，所以 \(CE^2=EB^2\). 由于线段长度非负，得到 \(CE=EB\).
\end{enumerate}
\end{solution}

\begin{problem}
如果已知 \(bx^2-4bx+2\left(a+c\right)=0\)（\(b\neq0\)）有两个相等的实数根，求证：\(a\)，\(b\)，\(c\) 成等差数列.
\end{problem}

\begin{solution}
因为二次方程有两个相等的实数根，所以其判别式等于零. 由 \(b\neq0\)，该方程确为二次方程，故
\[
\Delta=\left(-4b\right)^2-4b\cdot2\left(a+c\right)
=16b^2-8b\left(a+c\right)=0
\]
提取非零因子 \(8b\)，得 \(2b-a-c=0\)，即 \(a+c=2b\).
这正说明 \(b\) 是 \(a\)，\(c\) 的算术平均数，所以 \(a\)，\(b\)，\(c\) 成等差数列.
\end{solution}

\begin{problem}
\begin{enumerate}
\item 如图，为求河对岸某建筑物的高 \(AB\)，在地面上引一条基线 \(CD=a\)，测得 \(\angle ACB=\alpha\)，\(\angle BCD=\beta\)，\(\angle BDC=\gamma\)，求 \(AB\).

\item 如果 \(\alpha=30^\circ\)，\(\beta=75^\circ\)，\(\gamma=45^\circ\)，\(a=33\text{ 米}\)，求建筑物 \(AB\) 的高（保留一位小数）.

\end{enumerate}
\bitmapfigure[width=5.0cm]{img/1977/tianjin/q4_fig1.png}
\end{problem}

\begin{solution}
\begin{enumerate}
\item 建筑物的高 \(AB\) 垂直于地面，且 \(BC\) 在地面上，所以 \(AB\perp BC\). 在直角三角形 \(ABC\) 中，\(\tan\alpha=\dfrac{AB}{BC}\)，所以 \(AB=BC\tan\alpha\).

在三角形 \(BCD\) 中，
\(\angle CBD=180^\circ-\beta-\gamma\).
由正弦定理以及 \(CD=a\)，有
\[
\frac{BC}{\sin\gamma}
=\frac{CD}{\sin\angle CBD}
=\frac{a}{\sin\left(\beta+\gamma\right)}
\]
其中使用了 \(\sin\left(180^\circ-\theta\right)=\sin\theta\). 因此
\(BC=\frac{a\sin\gamma}{\sin\left(\beta+\gamma\right)}\).
代入 \(AB=BC\tan\alpha\)，得到
\(AB=\frac{a\sin\gamma\tan\alpha}{\sin\left(\beta+\gamma\right)}\).

\item 当 \(\alpha=30^\circ\)，\(\beta=75^\circ\)，\(\gamma=45^\circ\)，\(a=33\mathrm{~m}\) 时，
\[
BC=\frac{33\sin45^\circ}{\sin120^\circ}=11\sqrt{6}\mathrm{~m}
\]
\[
AB=BC\tan30^\circ
=11\sqrt{6}\cdot\frac{\sqrt{3}}{3}\mathrm{~m}
=11\sqrt{2}\mathrm{~m}
\]
因为 \(11\sqrt{2}\approx15.556\)，所以建筑物的高约为 \(15.6\mathrm{~m}\).
\end{enumerate}
\end{solution}

\begin{problem}
\begin{enumerate}
\item 求直线 \(3x-2y+1=0\) 和 \(x+3y+4=0\) 的交点坐标.

\item 求通过上述交点，并同直线 \(x+3y+4=0\) 垂直的直线方程.

\end{enumerate}
\end{problem}

\begin{solution}
\begin{enumerate}
\item 联立两条直线的方程，得
\[
\begin{cases}
3x-2y=-1,\\
x+3y=-4
\end{cases}
\]
将第二式乘以 \(3\)，得 \(3x+9y=-12\). 用此式减去第一式，得到 \(11y=-11\)，所以 \(y=-1\). 代入 \(x+3y=-4\)，得 \(x=-1\). 因此交点坐标为 \(\left(-1,-1\right)\).

\item 直线 \(x+3y+4=0\) 可写成 \(y=-\dfrac{1}{3}x-\dfrac{4}{3}\)，其斜率为 \(-\dfrac{1}{3}\). 与它垂直的直线斜率为 \(3\). 由点斜式，所求直线满足
\(y+1=3\left(x+1\right)\).
整理得所求直线方程为 \(y=3x+2\)，即 \(3x-y+2=0\).
\end{enumerate}
\end{solution}

\section{附加题}


\begin{problem}
求 \(\displaystyle\lim_{x\to0}\frac{\e^x-\e^{-x}-2x}{x-\sin nx}\) 的值.
\end{problem}

\begin{solution}
题中 \(n\) 视为固定参数，分 \(n\ne1\) 和 \(n=1\) 两种情况讨论.

当 \(n\ne1\) 时，分子、分母都趋于零，且分母的一阶导数在 \(x=0\) 处为 \(1-n\ne0\). 由洛必达法则，
\[
\lim_{x\to0}\frac{\e^x-\e^{-x}-2x}{x-\sin nx}
=\lim_{x\to0}\frac{\e^x+\e^{-x}-2}{1-n\cos nx}
=\frac{1+1-2}{1-n}=0
\]

当 \(n=1\) 时，连续使用洛必达法则，得到
\[
\begin{gathered}
\lim_{x\to0}\frac{\e^x-\e^{-x}-2x}{x-\sin x}
 =\lim_{x\to0}\frac{\e^x+\e^{-x}-2}{1-\cos x}
 =\lim_{x\to0}\frac{\e^x-\e^{-x}}{\sin x},\\
\lim_{x\to0}\frac{\e^x-\e^{-x}}{\sin x}
 =\lim_{x\to0}\frac{\e^x+\e^{-x}}{\cos x}=2
\end{gathered}
\]
因此原极限为
\[
\begin{cases}
0, & n\ne1,\\
2, & n=1.
\end{cases}
\]
\end{solution}

\begin{problem}
计算：\(\int_0^4\frac{x+2}{\sqrt{2x+1}}\mathrm{d}x\).
\end{problem}

\begin{solution}
令 \(t=\sqrt{2x+1}\). 则 \(t>0\)，且由 \(t^2=2x+1\) 得 \(x=\dfrac{t^2-1}{2}\)，\(\mathrm{d}x=t\,\mathrm{d}t\)，\(x+2=\dfrac{t^2+3}{2}\).
当 \(x=0\) 时 \(t=1\)；当 \(x=4\) 时 \(t=3\). 因此
\[
\int_0^4\frac{x+2}{\sqrt{2x+1}}\,\mathrm{d}x
=\frac{1}{2}\left[\frac{t^3}{3}+3t\right]_1^3
=\frac{1}{2}\left(18-\frac{10}{3}\right)=\frac{22}{3}
\]
\end{solution}
